\subsection{Certified enumeration and the six-pseudoline test}
\label{sec:counts}

\paragraph{C.01 (Proved): exhaustive encoding for a simple arrangement.}
Let $M$ be a simple, uniform oriented matroid of rank three on
$[n]=\{1,\ldots,n\}$, with a fixed chirotope $\chi_M$.
For a non-loop, non-parallel, same-rank extension by an element $p$, put
\[
 q_{ij}=\chi(i,j,p)\quad (1\leq i<j\leq n).
\]
The zero pairs of $q$ form a matching in the complete graph on $[n]$.
Conversely, we obtain every such extension by enumerating matchings,
assigning signs to the remaining pairs, and retaining precisely the
assignments satisfying the chirotope axioms on $[n]\cup\{p\}$.
We quotient by reorientation of $p$ by requiring the first nonzero
coordinate of $q$, in lexicographic pair order, to be positive.

\emph{Proof.}
If $q_{ij}=q_{ik}=0$ for distinct $i,j,k$, then $p$ lies in both
rank-two flats spanned by $\{i,j\}$ and $\{i,k\}$. Their union has
rank three because $M$ is uniform. Submodularity bounds the rank of
their intersection by one. That intersection contains $i$ and $p$,
so the non-loop $p$ is parallel to $i$, contrary to the hypothesis.
Thus the zero pairs are disjoint. Every extension therefore occurs in
our candidate list. We define the remaining triples by alternation,
set repeated triples to zero, and leave triples in $[n]$ equal to
$\chi_M$. The full signed basis-exchange axiom is then a necessary
and sufficient chirotope test, so accepting exactly its solutions
introduces no extra candidates. Multiplication of all $q_{ij}$ by
$-1$ is exactly reorientation of $p$. No allowed vector is zero, and
hence each orbit has exactly one representative with the stated
normalization. \hfill$\square$

For reproducibility, our fast enumerator first checks the necessary
three-term Grassmann--Pl\"ucker sign relations
\[
 \bigl\{
 \chi(a,b,c)\chi(a,d,e),
 -\chi(a,b,d)\chi(a,c,e),
 \chi(a,b,e)\chi(a,c,d)
 \bigr\}.
\]
Each such set must either be zero or contain both signs. This is
\emph{only a filter}: a second, independent program checks the full
signed basis-exchange axiom on every surviving assignment. Explicitly,
for every ordered basis $x=(x_1,x_2,x_3)$ and every ordered basis
$y=(y_1,y_2,y_3)$ it verifies that at least one $i$ satisfies
\[
 \chi(y_i,x_2,x_3)
 \chi(y_1,\ldots,y_{i-1},x_1,y_{i+1},\ldots,y_3)
 =\chi(x_1,x_2,x_3)\chi(y_1,y_2,y_3).
\]
It suffices to choose $x_1$ in each unordered basis, put the other two
entries in increasing order, and use increasing $y$: permutations of
the remaining entries multiply both sides by the same sign. We thus
check the full axiom, not merely its three-term consequence. For
$n=6$, the number of candidates, before this filtering, is
\[
 \sum_{k=0}^{3}
 \frac{6!}{(6-2k)!\,2^k k!}\,2^{15-k-1}=354304.
\]
The programs are \texttt{counts\_extensions.py}, its C++ inner loop
\texttt{counts\_extensions.cpp}, and
\texttt{counts\_independent\_check.py}. Their input and complete accepted
sign-vector lists accompany the paper.

\paragraph{C.02 (Proved): explicit input satisfying the test-case hypotheses.}
In homogeneous coordinates $[x:y:z]$ on $\mathbb{RP}^2$, take the six
lines with coefficient vectors
\[
\begin{split}
 a_1&=(1,0,1),\quad a_2=(2,0,1),\quad
 a_3=(0,1,1),\\
 a_4&=(0,3,1),\quad a_5=(1,1,1),\quad
 a_6=(4,4,1).
\end{split}
\]
Their twenty determinants, in lexicographic triple order, are
\[
\begin{split}
(&1,3,1,4,-2,-1,-7,-1,-13,-3,\\
 &-4,-1,-10,1,-14,-6,-2,-8,3,9).
\end{split}
\]
They are all nonzero; the arrangement is therefore simple and
nontrivial. The line $D=\{z=0\}$ contains
\[
 v_{12}=[0:1:0],\qquad
 v_{34}=[1:0:0],\qquad
 v_{56}=[1:-1:0].
\]
These are three distinct vertices. For $p=(0,0,1)$, its extension
vector in lexicographic pair order is
\[
 (0,+,+,+,+,+,+,+,+,0,-,-,-,-,0).
\]
All these claims follow by the displayed integer determinants or
cross products, which the attached programs also check.

\paragraph{C.03 (Proved): an independent disk-boundary certificate.}
We count the literal gen--gen inertia arcs of cote 154/120 independently
of the formula $s_1=2s_0$. Here are the complete rules used by the
checker, including the endpoint convention.

Realize the rank-two chirotope $q$ by oriented rays $w_i$ in an
oriented plane and order the $2n$ signed rays $\epsilon w_i$
cyclically. Equal signed rays form packets. For a two-element packet
$(\epsilon,i),(\eta,j)$, choose $k$ outside its unoriented parallel
class and order the first label before the second exactly when
\[
 -\eta\,\chi_M(i,j,k)q_{ik}=+1.                 \tag{C.1}
\]
The rank-two sign order and (C.1) suffice to construct the complete
boundary word. The global reversal of this word has no effect on
any count or gluing.

To verify (C.1), use local affine representatives
$a_i=(w_i,z_i)$ and $a_j=(\lambda w_i,z_j)$. At the endpoint in
direction $J(\epsilon w_i)$ of a sufficiently large circle,
the angular displacement has sign and relative order determined by
$\epsilon z_i/\lVert w_i\rVert$. Moreover
\[
 \det(a_i,a_j,a_k)
 =-(z_j-\lambda z_i)\det(w_i,w_k).
\]
For equal signed rays $\operatorname{sgn}\lambda=\epsilon\eta$;
substitution gives (C.1). This calculation is local to the
restriction on $\{i,j,k,p\}$ and does not assert a realization of
the whole extension. The two packets over the same vertex have
reversed orders, as required by the disk construction.

Each zero pair gives two generization intervals $G$, one in each
packet. For two adjacent endpoint labels $(\epsilon,i),(\eta,j)$
that belong to no $G$ packet, the interval is an ordinary
specialization interval $S$ precisely when, for every
$k\notin\{i,j\}$,
\[
 \chi_M(i,j,k)q_{ij}
     =\epsilon q_{ik}=\eta q_{jk}.             \tag{C.2}
\]
Indeed, the first expression is the side of line $k$ containing
$v_{ij}$ in the affine chart $p>0$. The other expressions are its
sides at the two endpoint rays. Agreement says that $k$ cuts neither
segment from $v_{ij}$ to those endpoints. Since two pseudolines
cross at most once and the boundary endpoints are adjacent, this
is equivalent to the intervening triangle being empty. Its
collapse is exactly the permitted specialization. The independent
checker confirms that the resulting pairs are exactly those
obtained by replacing one nonzero coordinate of $q$ by zero in
the exhaustive extension list.

For a possible exceptional sweep at $D_i$, the criterion is
\[
 q_{ab}\chi_M(a,b,i)=r\in\{+1,-1\}
 \quad\hbox{for every }a,b\ne i.              \tag{C.3}
\]
The common sign says that every base vertex off $D_i$ lies on the
same side of $D_i$ in the affine chart $p>0$. The other side is an
empty corridor. Its boundary interval lies on the side with sign
$-r$, and is shortened past the adjacent ordinary transition
intervals. Its position among the intervals is obtained from that
boundary subarc; the program uses its midpoint. When $q$ has one
zero pair, (C.3) is tested at the two lines through that vertex and
inserts the exceptional interval between the appropriate two
$G$ intervals. When $q$ has at least two zero pairs, there is no
such exceptional interval, since a zero pair avoiding $i$ would
violate (C.3).

These rules are exhaustive: the manuscript's ordinary transitions
are exactly the empty triangles and the parallel packets
(154/035--038), and its exceptional transitions are exactly the
empty corridors (154/036--038, together with the stable sweeps in
$\mathcal X_2$, 154/005). Consecutive $G$ intervals in the complete
cyclic list are therefore precisely the gen--gen inertia arcs.
Every boundary word and transition interval is recorded in the
boundary-certificate JSON files.

\paragraph{C.04 (Proved): the finite surface certificates.}
For the arrangement of C.02 the numbers $c_\nu$ of non-supercritical
relative positions through exactly $\nu$ vertices are
\[
 (c_0,c_1,c_2,c_3)=(80,150,75,10).
\]
In particular $s_2=315$. The direct boundary count of C.03 gives
$240$ gen--gen arcs at positions with two vertices and $60$ at
positions with three vertices. There are none at stable or
semi-stable positions. The $60$ flags of type $(0,1)$ in the
original simple arrangement contribute the exceptional vertices.
Thus the literal count in 154/120 is
\[
             s_0=240+60+60=360.               \tag{C.4}
\]

We also independently construct the complete edge identifications,
rather than infer the links from this count. An ordinary edge joins
an $S$ interval to the $G$ interval of its specialization with the
same unordered signed endpoint labels. The endpoint map exchanges
the two chord names. This is the reflection convention selected on
154/039, where one chord continues into the other chord; it is the
alternative to the discarded same-chord continuation. An exceptional
sweep at $i$ replaces $q$ by the normalized sign vector
\[
 \operatorname{sw}_i(q)_{ab}=
 \begin{cases}
   -q_{ab},&i\in\{a,b\},\\
   q_{ab},&i\notin\{a,b\}.
 \end{cases}                                  \tag{C.5}
\]
This is the local passage from one side of $D_i$ to the other.
For a stable exceptional edge, its endpoint map preserves the base
vertex at each end of the associated original edge. The two
adjacent $S$ intervals identify those vertices. For an unstable
exceptional edge, its endpoint map preserves the adjacent original
edge: resolve the adjacent $G$ interval to a stable position and
match its sweep partner on the other side. These are exactly the
four positions at an original vertex--edge flag.

The program \texttt{counts\_gluing.py} records every polygon, edge
pairing and endpoint permutation. A polygon corner is a consecutive
pair of transition arcs. Taking equivalence classes under the edge
endpoint maps yields $360$ classes. Every class contains exactly
four corners from four distinct polygons, and its link graph has
four distinct edges and degree two at each of its four vertices.
Thus every vertex link is a genuine four-cycle. Each edge is
paired once with an edge in a different polygon. No polygon has
two corners in the same equivalence class. Consequently each
closed polygon, edge and vertex embeds in the quotient; no
unverified closed-colimit embedding assertion is used. The finite
gluing relation is closed, so the quotient is compact Hausdorff.
The interior, edge, and vertex neighborhood tests prove that it is
a closed surface.

The adjacency graph of the polygons is connected. The signed
orientation constraints on adjacent polygons are inconsistent;
the full certificate contains an explicit orientation-reversing
cycle. These facts are checked directly from the endpoint maps.
There are $660$ ordinary edges and $60$ exceptional edges, hence
\[
 (s_0,s_1,s_2)=(360,720,315),\qquad
 \chi(\mathcal X)=-45.
\]
By the classification of compact connected surfaces,
$\mathcal X\cong\#^{47}\mathbb{RP}^2$ for this arrangement.
In particular the claim $\chi(\mathcal X)=1-n$ is false.

For clarity about the finite certification, every candidate is
specified by an explicit sign vector, every accepted candidate
satisfies the full chirotope axiom, and every gluing is an explicit
finite permutation of edge endpoints. The asserted counts and
link properties are exhaustive computations on those lists, with
assertions verifying each condition. The full axiom checker makes
$1\,089\,945$ signed basis-exchange checks for C.02. The computation
is therefore a complete finite counterexample, not a statistical
search or an inference from a partial enumeration.

\paragraph{C.05 (Proved): dependence on the arrangement at fixed $n$.}
As a second six-line input, take coefficient vectors
$(t,t^2,1)$ for $t=0,1,\ldots,5$. The same complete procedure gives
\[
 (c_0,c_1,c_2,c_3)=(74,132,57,4),\qquad
 (s_0,s_1,s_2)=(288,576,267).
\]
All $288$ vertex links are four-cycles, all closed faces embed,
and the quotient is connected and non-orientable. Thus this
surface is $\#^{23}\mathbb{RP}^2$, with Euler characteristic $-21$.
The independent axiom check makes $928\,233$ signed basis-exchange
checks. Together with C.04, this proves that the topology is not
a function of $n$ alone.

\begin{center}
\begin{tabular}{lrrrrrr}
\hline
Input & $c_0$ & $c_1$ & $c_2$ & $c_3$ & $s_0$ & $s_2$\\
\hline
Three lines & 4 & 6 & 0 & 0 & 12 & 10\\
Four convex dual points & 12 & 18 & 3 & 0 & 36 & 33\\
Five convex dual points & 31 & 50 & 15 & 0 & 100 & 96\\
Six convex dual points & 74 & 132 & 57 & 4 & 288 & 267\\
Six-line input C.02 & 80 & 150 & 75 & 10 & 360 & 315\\
\hline
\end{tabular}
\end{center}
The first three rows reproduce the earlier small cases and extend
the test to one five-line input. They are not substituted for a
claim about arbitrary arrangements.

\paragraph{C.06 (Proved): the fixed straight-line subproblem differs.}
For the particular realization C.02, the fifteen dual lines have
$42$ ordinary double vertices, one ordinary triple vertex at
$p=(0,0,1)$, and six supercritical vertices of multiplicity five.
Exact cross products give
\[
 V=49,\qquad E=42\cdot2+3+6\cdot5=117,\qquad
 F=1-V+E=69.
\]
The non-supercritical \emph{straight} positions therefore number
$(V-6)+E+F=229$, and their sector count is $2E=234$.
Their corresponding surface has Euler characteristic $-5$.
These counts are independently produced by
\texttt{counts\_straight\_dual.py} using integer arithmetic.
They do not equal the pseudoline counts $315$ and $360$.
Stretchability after moving the original lines cannot identify
all pseudoline extensions with straight extensions of a fixed
realization. In particular, around the displayed triple position
there are eight topological stable resolutions, while the fixed
straight realization supplies only six angular sectors.

\paragraph{C.07 (Proved): the first size allowing three-vertex positions.}
If a non-supercritical pseudoline $D$ passes through $\nu$ vertices
of a nontrivial arrangement of $n$ distinct pseudolines, then
$2\nu\le n$. Each such vertex uses at least two original lines;
no original line can occur at two different vertices on $D$,
because it would meet $D$ twice. The sets of original lines used
at those vertices are therefore disjoint. This proves the
inequality and shows that $n=6$ is the first possible size for
a three-vertex position. C.02 realizes this bound.
